% !TeX program = xelatex
\documentclass[12pt, a4paper, titlepage]{article}
\usepackage{xeCJK}
\usepackage{amsmath, amssymb}
\usepackage{geometry}
\usepackage{longtable}
\usepackage{booktabs}
\usepackage{tabularx}
\usepackage{caption}
\usepackage{silence}
\WarningFilter{caption}{Unknown document class}
\WarningFilter{xeCJK}{Redefining CJKfamily}
\WarningFilter{relsize}{Failed to get list}

% -- Fonts --
\setmainfont{Times New Roman}
\setCJKmainfont{Hiragino Mincho ProN}

\geometry{reset, margin=2.5cm}

\title{気液二相流の高精度数値解析における仮想時間反復と保存的フラックス・フィルタリングを統合した強結合スキームの理論的考察}
\author{}
\date{}

\begin{document}

\maketitle

\section*{1. 序論 (Introduction)}
Conservative Level Set (CLS) 法をはじめとする界面追跡法を用いた気液二相流シミュレーションにおいて、空間方向の高精度化（例：Compact Central Difference; CCD法による $O(h^6)$ 精度）は飛躍的な進歩を遂げている。しかし、時間方向の高精度化は依然として困難な課題である。単相流であれば高次Runge-Kutta (RK) 法などを素直に適用できるが、二相流においては界面を挟んだ密度および粘性の極端な不連続性（$\rho_l/\rho_g \sim 1000$）が誤差を増幅させる。加えて、非圧縮性流体の支配方程式は圧力Poisson方程式（PPE）を介した強い拘束条件を持つため、流体場全体に多段のRK法を適用することは、計算コストの観点から非現実的である。そのため、現状では界面移流にRK3法を用いつつ、流体の時間積分にはAdams-Bashforth 2次精度 (AB2) とIncremental Projection Method (IPC) を組み合わせた半陰的な手法が主流となっているが、これによる時間精度の頭打ちが空間高精度化の恩恵を制限している。本稿では、時間方向の高精度化を阻む根本的な要因を整理し、これを解決するための理論的枠組みとして、「仮想時間反復（Dual Time Stepping）」による流体・界面の強結合と、「保存的フラックス・フィルタリング」を統合した新しい数値計算スキームを提案し、その理論的妥当性を考察する。

\section*{2. 高次Projection法の限界と課題}
時間積分を高次化しようとする際、特にProjection法において以下の3つの主要な困難が顕在化する。
\begin{itemize}
    \item 界面物性の補間誤差: 界面近傍での物性値（密度・粘性）の急激な変化により、時間積分における実効精度が理論上の $\mathcal{O}(\Delta t^2)$ から $\mathcal{O}(\Delta t)$ へと退化する現象が不可避的に生じる。
    \item 変密度圧力Poisson方程式の計算負荷: 空間高精度化と整合する高次PPEは極めて条件数が悪化しやすい。時間積分の各中間ステップでこれを解くことは膨大な計算コストを要求する。
    \item 寄生流れ（Spurious Currents）の励起: 表面張力と圧力勾配の離散化演算子の不整合に起因する非物理的な渦流は、時間ステップを進めるごとに蓄積し、高次精度化による微小な振動（ギブス現象）と干渉して計算の破綻を招く。
\end{itemize}

\section*{3. 提案手法：強結合・保存的フィルタリングスキームの構築}
上述の課題を克服し、真の高精度化を達成するために、以下の2つのアプローチを統合したスキームを提案する。

\subsection*{3.1. 仮想時間反復による完全陰的結合（強結合）}
従来の分離解法（Operator Splitting）では、界面位置 $\phi$ の更新と速度場 $\mathbf{u}$ の計算が時間的にずれており、これが表面張力主導の流れにおいて致命的な誤差（Splitting Error）を生む。これを解決するため、各物理時間ステップ $n \to n+1$ の間に仮想時間 $\tau$ を導入し、反復インデックス $k$ を用いた Inner Iterations を実行する。
$$\frac{\mathbf{u}^{n+1, k+1} - \mathbf{u}^n}{\Delta t} + \mathcal{N}(\mathbf{u}^{n+1, k}, \phi^{n+1, k}) = -\frac{\nabla p^{n+1, k+1}}{\rho(\phi^{n+1, k})} + \frac{1}{\rho(\phi^{n+1, k})} \mathbf{F}_s(\phi^{n+1, k})$$反復が収束した極限 ($k \to \infty$) において、界面形状（曲率からの表面張力 $\mathbf{F}_s$）、移流項 $\mathcal{N}$、および圧力勾配が完全に時刻 $n+1$ でバランスする「フル・インプリシット（全陰解）」な状態が実現される。

\subsection*{3.2. 保存的フラックス・フィルタリング}
高次空間差分法（CCD法等）は優れたスペクトル分解能を持つが、不連続面近傍での高波数ノイズが避けられない。変数を直接平滑化すると物理量が散逸するため、「フラックス形式」に対してローパスフィルターを適用する。対象となる移流フラックス $F$ に対し、以下のようにフィルタリングされたフラックス $F^{Filtered}$ を定義する。$$F_{i+1/2}^{Filtered} = F_{i+1/2}^{CCD} - \mathcal{D}_{i+1/2}$$
ここで $\mathcal{D}$ は高波数成分のみを選択的に減衰させる高次の散逸項である。この形式は、隣接セル間でフラックスの授受が完全に相殺されるため、系全体の質量および運動量の厳密な保存を保証する。

\section*{4. 理論的考察と検証}
\subsection*{4.1. 幾何学的整合性と Balanced-Force 条件の回復}
仮想時間反復の最大の理論的利点は、「時間精度の幾何学的整合性」の獲得にある。分離解法では不可避であった「古い界面形状から計算された表面張力で新しい速度場を駆動する」という矛盾が解消される。収束状態では $\nabla p \approx \mathbf{F}_s$ が同一時刻で厳密に評価されるため、Balanced-Force条件が時間離散化のレベルで満たされ、寄生流れの発生源が根本的に絶たれる。

\subsection*{4.2. スペクトル解像度の維持と Modified Wavenumber}
保存的フラックス・フィルタリングの導入は、Modified Wavenumber（修正波数）の観点から極めて合理的である。フィルターの特性を高次多項式で設計することにより、物理的に意味のある低〜中波数域（マクロな流体構造や滑らかな界面変形）の振幅・位相特性を一切汚染することなく、計算格子で解像不可能なナイキスト波数付近の数値的エイリアシングのみをピンポイントで除去できる。これにより、高次スキームの「切れ味」と数値的安定性が両立する。

\subsection*{4.3. 保存則の数学的厳密性と反復収束性}
フラックス形式でのフィルタリングは、離散空間におけるガウスの発散定理を厳密に満たす。これにより、質量や運動量が機械精度（Machine Epsilon）の範囲内で保存される。この保存性の担保は、物理的妥当性のみならず、数値線形代数の観点からも極めて重要である。非保存的なエラーは仮想時間反復において「擬似的な湧き出し/吸い込み」として作用し、行列ソルバー（PPEなど）の条件数を悪化させるが、本手法ではこの阻害要因が排除されるため、Inner Iterations の収束性の大幅な向上が期待できる。

\section*{5. 結論と今後の展望 (Conclusion)}
本稿では、気液二相流シミュレーションの時間方向高精度化における課題を整理し、「仮想時間反復」と「保存的フラックス・フィルタリング」を統合した強結合スキームの理論的枠組みを提示した。本スキームは、Operator Splitting Errorの排除による時間的・幾何学的な整合性の確保と、高波数ノイズの選択的除去による厳密な保存則の維持を同時に達成するものであり、高次精度と安定性のトレードオフを打破する強力なアプローチである。今後の課題として、この強結合の枠組みにおいて、空間的に高精度に評価された界面曲率（$\kappa$）を、仮想時間反復の過程でいかに時間的に滑らか（$C^1$ 連続的）にフィードバックさせるかという、幾何学的更新アルゴリズムの更なる洗練が挙げられる。これが達成されれば、寄生流れを完全に抑制した極めてHigh-fidelityな二相流解析が実現すると確信される。

\end{document}